\section{De las máscaras marcadas al perfil}\label{app:transfer}
Dos casillas pueden depender de la supervivencia de una misma línea. Tratar sus indicadores
como independientes cambiaría la media del par, incluso si el perfil de cada casilla fuese
exacto. Por ello transferimos la contribución ponderada de cada estado con su propia masa.
Para un patrón fijo identificamos las etiquetas afines iguales y tomamos cada indicador
distinto una sola vez. Con \(d_*\) índices libres,
una coincidencia accidental satisface una ecuación afín no constante adicional y cuesta
\(O(n^{-1})\) al dividir por \(n^{d_*}\). En la trayectoria regular auxiliar, a \(x\ge b\), los
riesgos de las líneas axiales y diagonales vivas son \(1/m\) y \(\ell_f/m\), respectivamente.
Los segundos distan de sus valores deterministas a lo sumo \(\eta/n\) a precisión fija.
Dos líneas distintas se ocupan simultáneamente con probabilidad \(O_b(n^{-2})\). Bonferroni da
el producto de sus supervivencias de un paso con error \(O_{b,v}(\eta/n+n^{-2})\).
Multiplicamos cada indicador actual por su supervivencia determinista futura hasta el corte.
Una pieza altera \(O_v(n^{d_*-1})\) sumandos. Azuma para la media centrada, con saltos
\(O_{b,v}(1/n)\), y el control de regularidad dan una cola exponencial a precisión fija.
La unión sobre líneas marcadas es polinómica. Los cocientes de densidades hasta la parada
preservan la cola, incluida la cota \(\exp(O_b(n^{2/3}))\) para \(X\) al cambiar a \(\mu\).
Así, en la media normalizada cada etiqueta axial viva aporta \(b\) y cada etiqueta diagonal
distinta aporta \(U_f(b)\).

En particular, las sumas por línea viva de \(YW_b^0/j\) tienen márgenes \(1+o_{\Pr}(1)\)
uniformemente y dos columnas vivas tienen \(\ge c_bj\) vecinos comunes. En potenciales
logarítmicos de columna, el mapa que normaliza primero filas y después columnas tiene
derivada estocástica de entradas \(\ge c'_b/j\) en la bola de oscilación unidad. Su
contracción en oscilación da \(jD_{il}=Y_{il}W_b^0(i,l)(1+o_{\Pr}(1))\) para el escalado
de Sinkhorn, uniformemente en el soporte. Al retirar las dos casillas distinguidas, cada
margen pierde \(O_b(1/j)\) y dos columnas pierden a lo sumo dos vecinos comunes. La misma
prueba sigue aplicándose. La ocupación estacionaria de torres cumple
\(p_{il}/D_{il}=1+O_b(1/j)\) en el soporte. En los patrones integrados, las incidencias desde una
fila adicional \(y\) hacia las columnas marcadas tienen ley producto con parámetros
\(v_b(y,c_0),v_b(y,c_1)\). Las distintas columnas dan etiquetas diagonales distintas en cada
familia, salvo los \(O(1)\) valores de \(y\) que coinciden con una etiqueta distinguida.
Se puede multiplicar la media por cualquiera de los cuatro indicadores acotados de estado.
Para el término cuadrático de \(\eta\), sean \(B=p_b-1\),
\(G_R=BB^*, T_R=G_R(I-G_R)^{-1}, U=T_RB\), con los operadores transpuestos en columnas.
Sean \(d_R=\operatorname{diag}(T_R-G_R)\) y su análogo \(d_C\), con
\(\bar d=\int d_R=\int d_C\). Definimos
\(\psi_R=2[(B+U)d_C-T_Rd_R]-d_R\) y su versión de columnas. En los cuatro extremos
marcados (dos filas \(r_0,r_1\) y dos columnas \(c_0,c_1\)), las identidades de derivadas dan
\begin{align*}
 R''_1(0)/2={}&XW^0_{01}+YW^0_{10}+\sum_{\square} U-T_R(r_0,r_1)-T_C(c_0,c_1)
              -\tfrac12\sum_* d,\\
 A'_1(0)={}&-1+2\bar d+\tfrac12\sum_* \psi .
\end{align*}
El asterisco recorre los cuatro extremos y el cuadrado sus cuatro casillas; se ordenan las columnas
para que los casillas distinguidas sean \(00,11\). Para la primera identidad se derivan los
márgenes en el parámetro de ocupación. En los tipos ordinarios la respuesta media resuelve
\((I+N)y=N\alpha\), donde \(N\) tiene bloques cruzados \(p_b,p_b^*\) y \(\alpha\) inserta
los dos extremos de cada lado. Los bloques centrados de \(N(I+N)^{-1}N\) son \(T_R,-U,-U^*,T_C\).
En cada vértice marcado una incidencia repetida da su varianza de Bernoulli. Para el prefactor
se derivan el factor local, las diagonales de márgenes, la masa total y el determinante centrado.
Su suma es \(1+\langle d-1,y\rangle-\sum_*d/2\); sustituir la solución centrada y su modo
constante da la segunda identidad. La parte lineal límite \(jD\) es la suma \(L_s\)
en cada submedida. Al elevar al cuadrado resulta \(L_s^2/2\).

El argumento de patrones sirve a cualquier longitud fija, pero el segundo coeficiente contiene
la inversa de márgenes y, por tanto, recorridos de longitud arbitraria. Truncamos su expansión
en recorridos alternantes centrados. Cada recorrido de longitud \(q\ge2\) entre extremos fijos
está acotado por
\(C_b^2r_b^{q-2}\), por las dos normas de sección y la brecha \(r_b<1\). A orden dos en el
logaritmo del determinante hay a lo sumo dos inserciones en cada traza; su cola está acotada
por \(C_b\sum_{q>N}(1+q)^2r_b^{\max(2q-C,0)}\). A truncación fija, las recurrencias de márgenes
adjuntan árboles enraizados a una casilla o a una traza alternante. Las coincidencias de
etiquetas libres distintas cuestan \(O_{b,N}(1/n)\). Se conserva cada indicador repetido antes de
tomar la media. Una traza de longitud dos en una arista libre de media condicional \(V\) usa
\(Y^2=Y\) a la ida y vuelta. El exceso respecto al producto de medias aporta
\(V(1-V)\ell^2/(2d_i d_l)\), con \(\ell\) el peso normalizado de la arista en la Hessiana y
\(d_i,d_l\) sus grados de fila y columna. Las aristas raíz ya tienen fijados sus indicadores.
Así la varianza ordinaria se cuenta una vez. Primero \(n\to\infty\) a truncaciones fijas;
después se retiran las truncaciones. Sólo entonces \(b\downarrow0\): la positividad de \(v_b\)
da una minoración uniforme de \(p_b^*p_b\) cerca de cero y convergencia dominada en cada integral.
Con \(\mathcal I_0(C)=\sum_o\int_\Omega w_{o,0}\sum_s q_{o,s,0}C_{o,s}\), se obtiene
\[
 \lim_{b\downarrow0}\limsup_n\left|
 n\mathbb E_\mu[\textstyle\sum_{\{e,f\}}p_ep_f B_{\rm phys}(e,f);F_b]
 -\mathcal I_0(B_{2,s})\right|=0. \tag{F2}
\]
Así la frontera firmada por pares es \(J=-B_2+L_2\), con \(B_2=\mathcal I_0(B_{2,s})\)
y \(L_2=-\mathcal I_0(L_s^2/2)\). En el estado \(11\) cada orientación aporta su coeficiente
de una línea y la interacción de las dos orientaciones aparece una sola vez en \(-S\).
Sumar \(J\) a las otras dos medias da la constante \(P_T+J-S\) que entra en la comparación
de cortes, una vez tomado el límite en \(n\) para cada corte fijo.
